\documentclass[11pt,a4paper]{article}
\usepackage[margin=20mm,headheight=14pt,headsep=7mm,footskip=10mm]{geometry}
\usepackage[T1]{fontenc}
\usepackage[utf8]{inputenc}
\usepackage{lmodern}
\usepackage{microtype}
\usepackage{amsmath,amssymb}
\usepackage{booktabs,tabularx,array}
\usepackage{enumitem}
\usepackage{xcolor}
\usepackage{titlesec}
\usepackage{fancyhdr}
\usepackage{hyperref}
\urlstyle{same}

\definecolor{accent}{HTML}{174B60}
\definecolor{muted}{HTML}{536371}
\hypersetup{colorlinks=true,linkcolor=accent,urlcolor=accent,citecolor=accent,
  pdftitle={Predictive Risk-Guided Residual Control for Text-Driven Humanoid Manipulation},
  pdfauthor={Gu Wentao; He Yutong; Lu Yinuo; Shi Yixin; Sun Yuhan; Tie Yutong; Wang Qifa; Wang Yuxiang; Xia Fangxin; Yin Siyuan; Zang Ziyi; Zhu Jianyu}}
\setlength{\parindent}{0pt}
\setlength{\parskip}{4pt}
\setlist[itemize]{leftmargin=1.3em,topsep=3pt,itemsep=2pt,parsep=0pt}
\setlist[enumerate]{leftmargin=1.5em,topsep=3pt,itemsep=2pt,parsep=0pt}
\titleformat{\section}{\large\bfseries\color{accent}}{\thesection}{0.6em}{}
\titleformat{\subsection}{\normalsize\bfseries}{\thesubsection}{0.6em}{}
\titlespacing*{\section}{0pt}{10pt}{5pt}
\titlespacing*{\subsection}{0pt}{7pt}{3pt}
\pagestyle{fancy}
\fancyhf{}
\fancyhead[L]{\small\color{muted}Track 4 | Group 3}
\fancyhead[R]{\small\color{muted}Project Proposal}
\fancyfoot[C]{\small\thepage}
\renewcommand{\headrulewidth}{0.3pt}
\newcolumntype{Y}{>{\raggedright\arraybackslash}X}
\newcommand{\norm}[1]{\left\lVert#1\right\rVert}
\newcommand{\R}{\mathbb{R}}


\newcommand{\archnode}[6]{%
  \put(#1,#2){\color{#5}\rule{#3\unitlength}{#4\unitlength}}%
  \put(#1,#2){\color{accent!65}\framebox(#3,#4){%
    \color{black}\shortstack{#6}}}}
\raggedbottom

\begin{document}

\thispagestyle{fancy}
{\small\bfseries\color{accent}TRACK 4: HUMANOID-ROBOT MOTION GENERATION}\par
\vspace{4pt}
{\LARGE\bfseries Predictive Risk-Guided Residual Control\par
for Text-Driven Humanoid Manipulation\par}
\vspace{3pt}

\section{Topic, Scope, and Objectives}
We will study \textbf{text-driven tabletop block grasping with a simulated
Unitree G1}. Under the instructor-approved scope, we will compare one frozen
text-to-motion baseline with the same system augmented by \textbf{Predictive
Risk + Residual Correction}, using the same frozen SONIC controller.
All task data, evaluation, and demonstrations will come from simulation.

Our research question is: \textbf{Can a learned representation of where and
when execution problems may occur improve selective residual correction
beyond a binary gate or a scalar risk score?} Our objectives are to:
\begin{itemize}
  \item Establish a reproducible ARDY--SONIC grasp-and-lift baseline in simulation.
  \item Learn small, selective corrections while keeping ARDY and SONIC frozen.
  \item Isolate prediction, risk structure, and loss contributions; measure grasp
  success, response time, user corrections, and physical feasibility.
\end{itemize}
Residual robot control is established \cite{residual}; recent work also studies
gated body--hand residuals with frozen nominal modules \cite{gated}.
Our focus is a controlled study of \textbf{future-time/body-structured risk
conditioning} for text-driven contact tasks. Supervised correction will be
used in this project; reinforcement learning is outside the core scope.

\section{Selected Model and Simulation System}
\textbf{Model choice: ARDY-G1-RP-25FPS-Horizon8.}
ARDY supports streaming text and sparse spatial constraints, and its project
reports integration with SONIC \cite{ardy}. We will use the official
implementation \cite{ardycode} and G1 checkpoint \cite{ardycheckpoint}.
Its short chunks suit repeated correction. This selection reflects task fit;
its grasp success and hardware throughput must be measured.

\textbf{MuJoCo environment.}
We will combine SONIC \cite{sonic} and its official simulation setup
\cite{soniccode} with MuJoCo \cite{mujoco}, a free-base G1, table, dynamic block,
and articulated hand. Lifting must result from frictional hand contacts.
Standing constraints and SONIC provide lower-body stabilization without fixing
the base. Robot/hand assets, physics step, actuator gains, and contact settings
will be versioned and identical across comparisons. The initial task is
single-hand grasping of one block at reachable tabletop poses.

\textbf{Text grounding and execution.}
A shared sequencer will parse reaching, directional adjustment, grasping, and
lifting commands into phases and wrist goals. Current simulator object poses
provide geometric grounding; the core study explicitly uses privileged state.
The parser, constraints, and grasp logic will be documented, so language
performance is attributed to the complete pipeline. ARDY receives text,
pose history, and constraints; camera perception is deferred.

The initial B0 collection entry runs one episode by default or a sequential
batch with deterministic block XY positions. It uses a fixed right-hand
approach/settle/prepare/reach/lower/close/lift/hold protocol with text, measured history,
and sparse root-path/wrist/standing constraints supplied to frozen ARDY; frozen SONIC remains the
body controller. Immutable plans preserve seeds and completed failures on
resume. This offline collection mode pauses simulated time during generation;
it does not establish real-time operation or replace the planned instruction
suite. Pilot scene, hand, and contact thresholds must be calibrated before
freezing comparative evaluation. No teacher refinement participates in B0.
The implementation exposes two execution modes: batch collection and manual
JSON commands. Both start with an empty task; an explicit grasp flag adds the
tabletop task. Both modes load all frozen models before accepting or executing
prompts and explicitly report loading completion; planning-only runs load no
models. Manual commands are independent reset episodes, with an optional
persistent GUI; inputs received during generation or physical execution are
discarded. This interface is not yet the proposal's interactive-correction
subset. Sessions use timestamped batch/manual directories under
\texttt{output/}, using Asia/Shanghai \texttt{YYMMDD-HHMMSS} names with second
precision, with flat attempt folders and human-readable per-trial reports.
Rendering is a separate user-selected step after collection. Phase-focused
text is a pilot choice to compare against the previous full-task prompts with
matched seeds and scenes; it does not establish higher physical success.
By default, the tabletop task initializes the free root at the grounded table
target, then checks stability during a two-second stand before the five hand
phases. Walking is enabled only by an explicit \texttt{--walk} flag;
\texttt{--direct-start} remains a compatibility alias for the default start.
Manipulation-only and walk-and-grasp results are reported separately.
With \texttt{--walk}, the approach extension starts the root an additional
0.45 m behind a target
0.20 m before the front table edge, with configurable pilot distances. Actual
MuJoCo table/block/root state grounds the path without camera perception.
ARDY receives timed root waypoints with free feet for walking. Since the pinned
SONIC G1 interface does not consume world root XY, the sequencer checks actual
arrival rather than assuming generated root motion was tracked. Before reach,
require position error at most 6 cm, heading error at most 12 degrees, horizontal
speed at most .12 m/s, and both feet on the floor for .4 s continuously.
A missed or unstable arrival ends the trial. Rebuild standing anchors at the
arrived pose. A preparation phase then requests an 8-degree waist inclination
through ARDY torso-orientation conditions while holding the wrists back. The
pinned G1 has no articulated neck; this is a posture cue, not visual perception.
Repeat the continuous stability gate before reaching and anchor subsequent
torso goals to the measured prepared pose. These shared nominal torso conditions
do not expand the learned residual's arm-only output mask. Approach and
preparation time are included in the existing 30 s trial budget.
This change must be validated and frozen consistently across methods; it is a
new pilot protocol, not part of the earlier prompt-only comparison.
The current calibration sequencer holds the terminal walking reference during
settlement. Close and hold also preserve references instead of sampling another
motion. Reach/lower keep the fingers fully open by default; the historical
distal-finger preshape remains an explicit calibration control. Descent uses a
geometry-derived wrist-frame acquisition gate computed from the measured hand
center site and current block extents. If the gate is missed, bounded
lower-phase replanning requests a fresh ARDY reference from current execution
history. At acquisition, a checked measured-pose reference is held through
frozen SONIC while fingers close. Missing alignment after the bounded replans
or missing sustained opposing contact stops the trial. These shared state-based
nominal transitions are not learned residuals, teacher corrections, or extra
user instructions.
Both hands, including fingers and palms on the wrist-yaw bodies, may contact
the tabletop and table legs through normal contact physics. Other robot-table
contacts, non-right-hand robot-block contacts, non-foot robot-floor contacts,
and block-floor contacts remain prohibited. This revised pilot contact policy
is shared across methods; historical trials retain their original policy.
The success threshold remains a two-second hold, but collection continues
through the complete configured five-second hold phase, subject to failure
stops and the 30-second timeout. Later failures invalidate earlier success.
Threshold achievement is retained as a separate event; the block must also
remain above the clearance threshold with opposing contact at the recorded
episode end. A later loss onto the tabletop is a failed retained grasp even
without a prohibited collision. Historical success labels are not rewritten.
The shared pilot hand gains are now 6\,Nm/rad stiffness and
0.4\,Nm\,s/rad damping, under unchanged joint and motor torque bounds.
The shared tabletop contact solver uses Newton with elliptic friction cones,
impedance ratio 10 and tolerance $10^{-10}$, without NoSlip post-processing.
This calibrates soft-contact drift without changing mass, friction coefficients,
timestep, attachment policy, or the frozen models. These settings must be
identical across compared methods and frozen after validation.
The default generated lift lasts 4.8\,s; the historical 3.2\,s duration remains
configurable. Closure targets retain their previous values, while the acquisition
gate is derived from the measured hand-center site and block geometry. Longer
calibration holds and bounded alignment replans remain within the common 30-second
budget.
Two hand-only diagnostic trials have passed physical lift/contact/hold checks;
complete walk-to-grasp success is not yet established. Those partial successes
must not enter the complete-task success denominator or a verified expert set.

Use one high-level task instruction with automatic phase text, pose history,
and geometric constraints. ACT is not part of the core comparison: adding a
trained action-chunk policy would introduce a separate imitation-learning
system and demonstration budget. Bounded task-level state feedback and static
holds can be calibrated in the shared baseline without changing frozen models
or the learned arm-only correction interface.

The assignment comparison is between the nominal reference-generation pipeline
and learned risk-conditioned arm-reference correction through identical SONIC;
prompt changes alone do not define a second method. First calibrate the common
approach workspace, full hand clearance and grasp frame on validation data.
Do not weaken B0 to create a gain. Freeze common settings and perturbation levels
before held-out nominal, target-bias and action-latency comparisons. An arm-only
residual is not expected to repair walking failures or unreachable targets.
These calibration and learning milestones are not yet demonstrated successes.

For pilot assessment, ``held'' additionally requires opposing thumb and finger
normal contact above 0.01\,N throughout the two-second clearance interval;
this operational contact threshold remains subject to validation.

\textbf{State recording and offline visualization.}
Save the initial state and full MuJoCo robot, hand, and object state at every
50\,Hz control frame, together with a versioned compiled scene snapshot.
Offline rendering restores each saved state and recomputes derived geometry,
without stepping physics or running a policy. Use a fixed third-person camera
initially; the pilot tabletop collector also records body-attached head/wrist
cameras. The current visualization CLI defaults to 25 FPS RGB sampled from
50 Hz saved states, with explicit original state/frame mappings; 50 FPS is
available when all control frames are needed. Store RGB arrays
with camera names, simulation timestamps, and control-frame indices in
\texttt{vision/images.npz}, and optionally encode a video by simulation time.
Rendered images serve visualization and recording, not core inference inputs.
Visual replay snapshots do not replace the controller, buffer, and random-state
checkpoints required for counterfactual branches.

The adapter will use SONIC's joint-based streaming interface \cite{sonicstream}:
29 body-joint positions/velocities in the required order, body orientation,
and aligned frame indices. We will pin one matching encoder/decoder/config
set, initially the default release. ARDY output will be decoded, mapped, and
resampled into nominal references $A_t$:
\[
\text{Text + goals}\rightarrow\text{ARDY + adapter}\rightarrow A_t
\xrightarrow{\text{Risk + Residual}} A_t+\Delta A_t
\rightarrow\text{SONIC}\rightarrow\text{simulated G1}.
\]
Residuals change arm-joint references only; velocities and dependent reference
fields are recomputed afterward. Root, legs, and fingers receive no learned
offsets. A shared finger controller handles closure/release and contact checks.
All systems use identical reference limits, clearance checks, and failure stops.
A baseline bypasses the learned modules. Recovering falls through leg
corrections is outside this initial controller's authority.

\newpage
\section{System Architectures}
Both architectures share the same frozen models, hand controller, and physics.

\begin{figure}[!ht]
\centering
\setlength{\unitlength}{1mm}
\begin{picture}(170,64)
\linethickness{0.45pt}
\archnode{0}{36}{33}{18}{black!3}{Text + goals\\+ pose history}
\archnode{43}{36}{33}{18}{accent!8}{\textbf{ARDY} (frozen)\\Adapter + checks}
\archnode{86}{36}{33}{18}{accent!8}{\textbf{SONIC} (frozen)\\Body controller}
\archnode{129}{36}{33}{18}{accent!4}{\textbf{G1 in MuJoCo}\\Contact physics}
\archnode{43}{9}{42}{15}{black!3}{Shared finger controller\\Grasp / release}
\color{accent}
\put(33,45){\vector(1,0){10}}
\put(76,45){\vector(1,0){10}}
\put(119,45){\vector(1,0){10}}
\put(81,59){\makebox(0,0){\small Nominal reference}}
\put(16.5,36){\line(0,-1){19.5}}
\put(16.5,16.5){\vector(1,0){26.5}}
\put(85,16.5){\line(1,0){60.5}}
\put(145.5,16.5){\vector(0,1){19.5}}
\put(108,20){\makebox(0,0){\small Finger commands}}
\color{muted}
\put(162,45){\line(1,0){6}}
\put(168,45){\line(0,-1){43}}
\put(168,2){\line(-1,0){160}}
\put(8,2){\vector(0,1){34}}
\put(121,6){\makebox(0,0){\small Simulator feedback}}
\end{picture}
\caption{\textbf{Baseline (B0): ARDY + SONIC.} Frozen ARDY supplies nominal
references to SONIC. The shared task interface controls the fingers and
consumes simulator feedback.}
\label{fig:baseline}
\end{figure}

\begin{figure}[!ht]
\centering
\setlength{\unitlength}{1mm}
\begin{picture}(170,124)
\linethickness{0.45pt}
\archnode{0}{103}{33}{16}{black!3}{Text + goals\\+ pose history}
\archnode{43}{103}{33}{16}{accent!8}{\textbf{ARDY} (frozen)\\Motion adapter}
\archnode{88}{103}{37}{16}{black!3}{$A+\Delta A$\\Reference checks}
\archnode{136}{103}{30}{16}{accent!8}{\textbf{SONIC}\\Frozen}
\archnode{0}{63}{33}{18}{black!3}{Execution history\\$S$}
\archnode{43}{63}{33}{18}{orange!13}{\textbf{Risk model}\\$G_\phi(S,A)$}
\archnode{136}{63}{30}{18}{accent!4}{\textbf{G1 in MuJoCo}\\Contact physics}
\archnode{43}{22}{33}{20}{violet!7}{Structured risk\\tokens $Z$\\Gate score $p$}
\archnode{88}{22}{37}{20}{orange!13}{\textbf{Residual model}\\$R_\theta(S,A,Z)$\\Active if $p\geq\tau$}
\archnode{136}{22}{30}{20}{black!3}{Task phase\\Finger controller\\Shared}
\color{accent}
% Nominal path stays intact; it also conditions both learned modules.
\put(33,111){\vector(1,0){10}}
\put(76,111){\vector(1,0){12}}
\put(125,111){\vector(1,0){11}}
\put(151,103){\vector(0,-1){22}}
\put(81,122){\makebox(0,0){$A$}}
\put(82,111){\line(0,-1){59}}
\put(82,90){\line(-1,0){22.5}}
\put(59.5,90){\vector(0,-1){9}}
\put(82,52){\line(1,0){16.5}}
\put(98.5,52){\vector(0,-1){10}}
% History enters risk and residual separately.
\put(33,72){\vector(1,0){10}}
\put(37,72){\line(0,-1){61}}
\put(37,11){\line(1,0){67}}
\put(104,11){\vector(0,1){11}}
\put(72,14.5){\makebox(0,0){\small History context}}
% Full risk representation is passed to the residual, not just the gate.
\put(59.5,63){\vector(0,-1){21}}
\put(76,32){\vector(1,0){12}}
\put(82,37){\makebox(0,0){$Z,p$}}
\put(118,42){\vector(0,1){61}}
\put(109,84){\makebox(0,0){$\Delta A^{\mathrm{des}}$}}
\put(151,42){\vector(0,1){21}}
\put(157,52){\makebox(0,0){\small Fingers}}
% Simulator feedback closes the learned control loop.
\color{muted}
\put(166,72){\line(1,0){3}}
\put(169,72){\line(0,-1){68}}
\put(169,4){\line(-1,0){152.5}}
\put(16.5,4){\vector(0,1){59}}
\put(132,8){\makebox(0,0){\small Simulator feedback}}
\end{picture}
\caption{\textbf{Proposed (P): ARDY + Risk + Residual + SONIC.}
Risk uses history and future nominal motion. Structured tokens condition the
bounded residual; low risk requests zero correction. Orange modules are trainable;
blue modules are frozen. Reference checks smooth gate transitions; shared
hand control and feedback match Figure~\ref{fig:baseline}.}
\label{fig:proposed}
\end{figure}

\clearpage
\section{Predictive Risk and Residual Learning}
\subsection{Anticipating execution problems}
The risk model receives 200--500\,ms of execution history and future nominal
references, with common task-phase and planned finger-command context:
\[
S_t=[s_{t-K+1},\ldots,s_t],\qquad A_t=[a_t,\ldots,a_{t+H-1}],
\qquad (Z_t,p_t)=G_\phi(S_t,A_t).
\]
State includes joint positions/velocities, base orientation/angular velocity,
hand--object pose, contacts, tracking error, and task phase. Current/past
observations are inputs; future actual execution is used only for labels.
The hand-command context is observed but never changed by the residual.

A compact Transformer (width 256, four layers, eight heads) will produce
$Z_t\in\R^{H\times B\times D}$ with body/time positional embeddings:
$H=8$ samples spaced 40\,ms apart, $B=4$ groups (left/right arm-hand, torso,
lower body), and $D=32$. Body/time auxiliary heads supervise tracking error,
phase-dependent contact loss, and instability; a pooled head gives $p_t$:
\[
\mathcal L_{\mathrm{risk}}=
\lambda_e\mathcal L_{\mathrm{track}}+
\lambda_c\mathcal L_{\mathrm{contact}}+
\lambda_b\mathcal L_{\mathrm{balance}}+
\lambda_i\mathcal L_{\mathrm{intervene}}.
\]
Intervention labels mark threshold violations during nominal future execution;
thresholds are validated then fixed. Open-hand phases are not grasp failures.
Losses mask unavailable future steps. High risk indicates possible failure,
not a guarantee that an arm correction can recover it.

\subsection{Bounded correction and distinct supervision}
A temporal Transformer reads $(S_t,A_t,Z_t)$. An arm-only mask $M$ and bounds
$\epsilon$ restrict the prediction; gating selects a desired correction:
\[
\widehat{\Delta A}_t=M\odot\operatorname{clip}
\bigl(R_\theta(S_t,A_t,Z_t),-\epsilon,\epsilon\bigr),\qquad
\Delta A_t^{\mathrm{des}}=\mathbf{1}[p_t\geq\tau]\widehat{\Delta A}_t.
\]
The shared reference checks apply magnitude/rate limits before execution.
Low risk stops residual inference and requests zero; any previous offset
ramps to zero under the same rate limit. Diagrams use $\Delta A$ for the
applied offset after this step.

Teacher recoveries supply $\Delta A_t^*=A_t^{\mathrm{expert}}-A_t$.
Correction and identity losses use \textbf{disjoint} recoverable-perturbed and
stable-nominal sets $\mathcal D_c$ and $\mathcal D_s$:
\[
\begin{aligned}
\mathcal L_{\mathrm{res}}={}&
\mathbb E_{\mathcal D_c}\norm{\widehat{\Delta A}_t-\Delta A_t^*}_F^2
+\alpha\mathbb E_{\mathcal D_s}\norm{\widehat{\Delta A}_t}_F^2\\
&+\beta\mathbb E_{\mathcal D_c\cup\mathcal D_s}
\sum_{h=0}^{H-2}\norm{\widehat{\Delta a}_{t,h+1}-\widehat{\Delta a}_{t,h}}_2^2.
\end{aligned}
\]
This avoids counting stable zero targets twice. Sampling ratios and loss
normalization are fixed across ablations; removing identity supervision now
has a distinct meaning. Training uses bounded predictions before the gate.

\subsection{Simulation data and runtime}
An offline constrained IK teacher will propose arm-only corrections within
the same bounds, verified through frozen SONIC and the shared hand controller
from the \textbf{same perturbed state}. Simulated teleoperation may supplement
recoveries. Failed recoveries supervise risk only. Nominal-label branches
restore simulator state, controller history, reference buffers, and disturbance
seeds; corrected outcomes never substitute for nominal counterfactual labels.

Split parent episodes (including all branches) before window extraction. Train risk first,
then residuals with \textbf{predicted} tokens and a frozen risk encoder on S4000.
The teacher is absent at evaluation. Joint fine-tuning is deferred.

SONIC targets 50\,Hz; risk/residual updates target 10--20\,Hz. A timestamped
buffer covers both risk and SONIC lookahead; references are resampled to each
clock. Execute only until the next update before replanning. Underruns trigger
a shared hold and are logged. Simulation may run slower than wall time;
25\,FPS motion sampling does not establish real-time inference.

\newpage
\section{Primary Comparison and Ablation Experiments}
Every configuration uses the \textbf{same frozen ARDY and SONIC checkpoints},
scene, task sequencer, hand controller, observations, and reference adapter.
The primary comparison is \textbf{B0 versus P}; intermediate controls explain
where any improvement comes from.

\begingroup
\small
\renewcommand{\arraystretch}{1.13}
\begin{tabularx}{\linewidth}{@{}p{9mm}p{45mm}Y@{}}
\toprule
\textbf{ID} & \textbf{Configuration} & \textbf{Purpose} \\
\midrule
B0 & ARDY + SONIC & Frozen baseline; no learned correction. \\
B1 & Always-on residual & Residual receives history and nominal actions only; tests whether
ordinary residual learning is sufficient. \\
B2 & Reactive gate + residual & Current tracking/contact/balance thresholds
gate the B1 residual; tests reactive intervention. \\
I1 & Predictive gate only & Predicted risk activates a residual with no
risk features; tests anticipation without risk conditioning. \\
I2 & Predictive scalar + residual & Residual also receives the risk probability;
tests a single predictive risk score. \\
I3 & Body-wise risk + residual & Residual receives body-wise auxiliary risks;
tests spatial localization. \\
I4 & Body--time risk map + residual & Residual receives explicit auxiliary
predictions over $H\times B$; tests temporal localization. \\
P & Structured tokens + residual & Full proposed method: gated residual
conditioned on learned $H\times B\times D$ risk features. \\
\bottomrule
\end{tabularx}
\endgroup

\subsection{Isolating the risk interface}
I1--I4 and P share a frozen risk encoder and gate. B1/B2/I1 reuse one residual
checkpoint. Identical body/time slots contain zeros (I1), repeated probability
(I2), time-pooled body-wise auxiliary risks (I3), body/time predictions (I4),
or learned tokens (P). Threshold-normalized auxiliary channels share feature
width, decoder capacity, corrective data, and training budgets. We will report
adapter/decoder parameter counts.

The gate function is fixed, but closed-loop states can change activation rates.
Compare success--intervention curves at validation-selected thresholds and
residual errors on identical held-out windows. Structured features should
improve success/locality at comparable intervention effort; this is a hypothesis.

\subsection{Component and loss ablations}
Each row below changes \textbf{one factor} relative to P. Other losses, correction
bounds, data, and optimization budgets remain fixed.

\begingroup
\small
\renewcommand{\arraystretch}{1.1}
\begin{tabularx}{\linewidth}{@{}p{40mm}YY@{}}
\toprule
\textbf{Ablation} & \textbf{Controlled change} & \textbf{Question / evidence} \\
\midrule
Pooled learned tokens & Average body/time tokens and repeat over slots; retrain residual. &
Does localization matter beyond a generic learned feature? \\
No future-action input & Remove future actions from risk only; retrain risk/residual. &
Does anticipation improve warning lead time and success? \\
No execution history & Risk uses current state and future actions; retrain risk/residual. &
Does temporal context expose delay and drift? \\
No auxiliary risk losses & Train risk with intervention loss only; retrain residual. &
Do execution-related heads make tokens useful? \\
No selective gate & Force the trained P gate on; keep both networks fixed. &
Does selection reduce unnecessary corrections? \\
No identity loss & Retrain residual with $\alpha=0$. &
Does removing stable-state supervision increase unnecessary correction? \\
No smoothness loss & Retrain residual with $\beta=0$. &
Do reference jitter, contacts, or grasp stability worsen? \\
\bottomrule
\end{tabularx}
\endgroup

Hard bounds remain; compare raw and applied offsets because limiting can mask
loss effects. Prioritize B0/B1/I2/P and report any unfinished ablations.

\newpage
\section{Evaluation Protocol and Analysis}
\textbf{Prioritized test suite.}
Core conditions are nominal execution, grasp-target bias, and action latency.
The instruction set covers reaching, adjustment, grasping, lifting, paraphrases,
and left/right targets. Fixed pilot-calibrated perturbation levels will be
published before testing. Mass, friction, pushes, and held-out combinations
form a secondary stress suite after core comparisons are complete. Prompt
paraphrases and scene seeds are separated from training where applicable.

\textbf{Success and trial budget.}
The instructed block's lowest point must clear the tabletop by 5\,cm and remain
held for 2\,s within 30\,s of simulated time, without a fall or prohibited
collision. Fall/clearance thresholds will be fixed in the evaluation config.
Hand-table contact is permitted under the shared pilot contact policy above.
Report threshold achievement and final retained success separately; a later
drop before recording ends does not count as retained task success.
We plan \textbf{20 paired episodes per core condition per configuration}.
Match instructions, initial states, perturbations, and ARDY random seeds;
closed-loop nominal motions may subsequently diverge. B1, I2, and P target
three training seeds; other ablations start with one and are exploratory.
The full plan on core conditions is approximately 1,260 episodes (10.5 simulated
hours at the timeout limit), before secondary stress tests or runtime overhead.

\textbf{Metrics and uncertainty.}
Report success with Wilson 95\% intervals, paired success differences with
an episode-level bootstrap, and results per training seed. Overlapping windows
are not independent trials. Also report command-to-first-motion/completion
time, wrist error, slips/falls/collisions, intervention rate, correction norm,
locality, and smoothness. Risk metrics include precision--recall, false alarms,
and warning lead time on held-out nominal rollouts, avoiding successful
interventions being relabeled as false alarms. Latency includes text encoding,
generation, communication, risk, and residual inference (median/p95).
Simulated task time, wall-clock response, and simulation speed are separate.

Autonomous trials receive no extra user corrections. An interactive subset
allows at most two corrections under a fixed rule, excluding planned phase
commands from the count. Failed trials and timeouts remain in the denominator;
completion times are reported alongside success. Videos will cover language,
pose, contact, balance, intervention, and buffering failures.

\section{Milestones and Expected Deliverables}
Project period: \textbf{29 September--15 October 2026}.
Integration, data tooling, and model development will run in parallel.

\begingroup
\small
\renewcommand{\arraystretch}{1.12}
\begin{tabularx}{\linewidth}{@{}p{30mm}Y@{}}
\toprule
\textbf{Period} & \textbf{Milestone and acceptance evidence} \\
\midrule
Sep 29--30 & Verify model access and S4000 operators; launch G1/SONIC
in MuJoCo. Build data and evaluation tools in parallel. \\
Oct 1--3 & Integrate ARDY, constraints, and hand control; freeze the
baseline after a verified contact grasp/lift. Validate teacher recovery targets. \\
Oct 4--7 & Expand perturbation/correction data; train risk and residual
models. Validate selective closed-loop correction. \\
Oct 8--11 & Complete B0/B1/I2/P first; then expand interfaces and
component ablations. Freeze the completed result set and failures. \\
Oct 12--14 & Freeze results; complete report, slides, repository, and
videos. Reproduce key runs and rehearse the presentation. \\
Oct 15 & Final consistency checks and submission of all required
deliverables by the course deadline. \\
\bottomrule
\end{tabularx}
\endgroup

By Sep.~30, verify S4000 training operators, text-encoder access, and a working
SONIC inference backend \cite{ardycode}, \cite{soniccode}. TensorRT compatibility
is not assumed; confirm any separate compatible inference host is available.
If integration slips, reduce scene variation before expanding experiments.
No throughput or successful grasp is claimed before these checks.

Deliverables: the simulation system, checkpoints, ablation results, and a
public GitHub repository with setup/data/training/evaluation scripts; a
5--8-page report (excluding references); and a 10-minute presentation with
5-minute Q\&A. Videos show commands and successful/failed trials for both systems.

\newpage
\section{Future Directions: VLA Policies in Simulation}
After the ARDY study, we propose testing \textbf{GR00T N1.7} \cite{groot} and
\textbf{$\pi_{0.5}$} \cite{openpi} with the same risk-guided residual principle,
using rendered images, language, and robot state. GR00T documents a G1--SONIC
latent-action interface; $\pi_{0.5}$ would require G1-specific adaptation.
Each future baseline must first be validated and then frozen. Risk/residual
modules would be trained on that policy's simulated rollouts and evaluated
with and without correction under matched perturbations. Correction coordinates
and bounds must suit the action interface; adding offsets to latent tokens
is not automatically equivalent to correcting joint references. This is future
work, and zero-shot weight transfer would be a separate experiment.

\section{Team Roles and Responsibilities}
Each member owns a technical component and contributes its methods, evidence,
and limitations to the report and presentation.

\begingroup
\small
\renewcommand{\arraystretch}{1.1}
\begin{tabularx}{\linewidth}{@{}>{\raggedright\arraybackslash}p{26mm}>{\raggedright\arraybackslash}p{35mm}Y@{}}
\toprule
\textbf{Member} & \textbf{Primary role} & \textbf{Responsibilities} \\
\midrule
Gu Wentao & Architecture, experiments, and training lead &
Design the framework and experiments/ablations; coordinate module interfaces;
lead S4000 training, tuning, and checkpoints. \\
\addlinespace[2pt]
He Yutong & ARDY integration &
Integrate text encoding, pose/constraint inputs, and streaming inference;
validate the selected checkpoint. \\
\addlinespace[2pt]
Lu Yinuo & SONIC and adaptation &
Implement joint/frame mapping, resampling, and reference updates;
validate lower-body tracking and balance. \\
\addlinespace[2pt]
Shi Yixin & Simulation environment &
Build the MuJoCo G1/table/block scene, contact physics, cameras,
and reproducible resets. \\
\addlinespace[2pt]
Sun Yuhan & Predictive risk model &
Implement history encoding, auxiliary heads, structured risk tokens,
and prediction validation. \\
\addlinespace[2pt]
Tie Yutong & Residual correction &
Implement risk-conditioned residuals, correction bounds,
identity/smoothness losses, and model ablations. \\
\addlinespace[2pt]
Wang Qifa & Data engineering &
Build synchronized episode recording, loaders, normalization,
split audits, and optional LeRobot export. \\
\addlinespace[2pt]
Wang Yuxiang & Perturbations, labels, and mathematical reasoning  &
Generate perturbations, nominal snapshot rollouts, risk labels,
and validated corrective targets; analyze mathematical assumptions. \\
\addlinespace[2pt]
Xia Fangxin & Language grounding and goals &
Build the task sequencer, instruction suite, wrist/standing constraints,
and IK teacher interface. \\
\addlinespace[2pt]
Yin Siyuan & Evaluation and statistics &
Run paired trials and ablation comparisons; analyze uncertainty,
intervention trade-offs, and failure categories. \\
\addlinespace[2pt]
Zang Ziyi & Simulated hand interaction &
Implement finger actuation, contact checks, grasp/lift verification,
and hand-model validation. \\
\addlinespace[2pt]
Zhu Jianyu & Runtime and reproducibility &
Implement buffering, timing, launch scripts, video recording,
and the reproducible release. \\
\bottomrule
\end{tabularx}
\endgroup

{\small External code, models, and datasets will be acknowledged. LLM assistance
was used to draft this proposal; the team will review it under the course policy.\par}

\clearpage
\begingroup
\small
\raggedright
\setlength{\parskip}{1pt}
\begin{thebibliography}{11}
\setlength{\itemsep}{7pt}
\bibitem{residual}
T. Johannink et al., ``Residual reinforcement learning for robot
control,'' in Proc. IEEE Int. Conf. Robot. Autom. (ICRA), Montreal,
QC, Canada, 2019, pp. 6023--6029, doi:
\href{https://doi.org/10.1109/ICRA.2019.8794127}{10.1109/ICRA.2019.8794127}.

\bibitem{gated}
R. Wu, S. Li, L. Zhang, A. Knoll, and Z. Chen, ``Gated residual body-hand
coordination for whole-body humanoid teleoperation,'' 2026,
arXiv:2609.18763. [Online]. Available: \url{https://arxiv.org/abs/2609.18763}

\bibitem{ardy}
K. Zhao, M. Petrovich, H. Zhang, T. Wang, S. Tang, and D. Rempe,
``ARDY: Autoregressive diffusion with hybrid representation for interactive
human motion generation,'' ACM Trans. Graph., vol. 45, no. 4,
2026, Art. no. 86, doi:
\href{https://doi.org/10.1145/3811284}{10.1145/3811284}.

\bibitem{ardycode}
NVIDIA. ``ARDY.'' GitHub. Accessed: Sep. 28, 2026. [Online]. Available:
\url{https://github.com/nv-tlabs/ardy}

\bibitem{ardycheckpoint}
NVIDIA. ``ARDY-G1-RP-25FPS-Horizon8.'' Hugging Face.
Accessed: Sep. 28, 2026. [Online]. Available:
\url{https://huggingface.co/nvidia/ARDY-G1-RP-25FPS-Horizon8}

\bibitem{sonic}
Z. Luo et al., ``SONIC: Supersizing motion tracking for natural
humanoid whole-body control,'' Sci. Robot., vol. 11, no. 117,
2026, Art. no. eaed4592, doi:
\href{https://doi.org/10.1126/scirobotics.aed4592}{10.1126/scirobotics.aed4592}.

\bibitem{soniccode}
NVIDIA. ``GR00T-WholeBodyControl.'' GitHub.
Accessed: Sep. 28, 2026. [Online]. Available:
\url{https://github.com/NVlabs/GR00T-WholeBodyControl}

\bibitem{mujoco}
E. Todorov, T. Erez, and Y. Tassa, ``MuJoCo: A physics engine for model-based
control,'' in Proc. IEEE/RSJ Int. Conf. Intell. Robots Syst. (IROS),
Vilamoura, Portugal, 2012, pp. 5026--5033, doi:
\href{https://doi.org/10.1109/IROS.2012.6386109}{10.1109/IROS.2012.6386109}.

\bibitem{sonicstream}
NVIDIA. ``Streaming motion tracking.'' GR00T-WholeBodyControl
Documentation. Accessed: Sep. 28, 2026. [Online]. Available:
\url{https://nvlabs.github.io/GR00T-WholeBodyControl/tutorials/zmq.html}

\bibitem{groot}
NVIDIA. ``Isaac-GR00T N1.7.'' GitHub.
Accessed: Sep. 28, 2026. [Online]. Available:
\url{https://github.com/NVIDIA/Isaac-GR00T}

\bibitem{openpi}
Physical Intelligence. ``OpenPI.'' GitHub.
Accessed: Sep. 28, 2026. [Online]. Available:
\url{https://github.com/Physical-Intelligence/openpi}
\end{thebibliography}
\endgroup
\end{document}
